% Part: lambda-calculus
% Chapter: introduction
% Section: church-rosser

\documentclass[../../../include/open-logic-section]{subfiles}

\begin{document}

\olfileid{lam}{int}{cr}
\olsection{ویژگی چرچ--راسر}

\begin{thm}
\ollabel{thm:church-rosser}
فرض کنید~$M$، $N_1$ و~$N_2$ ترم‌هایی باشند چنان‌که~$M \red N_1$ و
$M \red N_2$. آنگاه ترمی چون~$P$ وجود دارد که~$N_1 \red P$ و
$N_2 \red P$.
\end{thm}

\begin{cor}
فرض کنید~$M$ را بتوان به صورت نرمال کاهش داد. آنگاه این صورت نرمال
یکتاست.
\end{cor}

\begin{proof}
اگر~$M \red N_1$ و~$M \red N_2$، بنا بر قضیهٔ پیشین ترمی چون~$P$ وجود
دارد که~$N_1$ و~$N_2$ هر دو به~$P$ کاهش می‌یابند. اگر~$N_1$ و~$N_2$ هر
دو در صورت نرمال باشند، این امر تنها هنگامی ممکن است که
$N_1 \ident P \ident N_2$.
\end{proof}

سرانجام، می‌گوییم دو ترم~$M$ و~$N$ \emph{$\beta$-هم‌ارز} یا صرفاً
\emph{هم‌ارز} هستند هرگاه به ترمی مشترک کاهش یابند؛ به بیان دیگر، هرگاه
$P$ای وجود داشته باشد که~$M \red P$ و~$N \red P$. این را به‌صورت
$M \equal[\beta] N$ می‌نویسیم. با استفاده از
\olref{thm:church-rosser} می‌توانید بررسی کنید که~$\equal[\beta]$ رابطه‌ای
هم‌ارزی است و ویژگی افزودهٔ زیر را دارد: برای هر~$M$ و~$N$، اگر
$M \red N$ یا~$N \red M$، آنگاه~$M \equal[\beta] N$. (در واقع می‌توان
نشان داد که~$\equal[\beta]$ \emph{کوچک‌ترین} رابطهٔ هم‌ارزیِ دارای این
ویژگی است.)

\end{document}

